\chapter{Threat Model}
\label{ch:threat}

A defence is only meaningful relative to an explicit statement of what the attacker can do and
what the defender is trying to preserve. This chapter gives that statement for
\textsc{EdgeGuard}. \Cref{sec:tm:system} describes the system and its trust boundaries,
\cref{sec:tm:adversary} the adversary, \cref{sec:tm:attacks} the attack classes considered, and
\cref{sec:tm:goals} the properties \textsc{EdgeGuard} aims to provide and those it does not.

% ---------------------------------------------------------------------------
\section{System model}
\label{sec:tm:system}

\begin{figure}[t]
\centering
\begin{tikzpicture}[
  comp/.style={draw=eggray, rounded corners=2pt, align=center, font=\small, minimum height=0.9cm, inner sep=4pt},
  trusted/.style={comp, draw=egblue, fill=eglight},
  atk/.style={draw=egorange, fill=egorange!10, rounded corners=2pt, align=center, font=\scriptsize, inner sep=3pt},
  arr/.style={-{Stealth}, eggray},
  aarr/.style={-{Stealth}, egorange, thick, dashed}]
\node[comp, minimum width=2.6cm] (vendor) at (0,0) {Model vendor\\(signs weights)};
\node[comp, minimum width=2.6cm] (repo) at (4.2,0) {Distribution\\channel / repo};
\node[comp, minimum width=3.2cm] (dram) at (4.2,-2.6) {DRAM\\weights $\{L_i\}$};
\node[comp, minimum width=2.6cm] (npu) at (8.6,-2.6) {NPU (HTP)\\inference};
\node[trusted, minimum width=3.4cm] (eg) at (8.6,-5.2) {\textsc{EdgeGuard} verifier\\+ public key + calib.\ stats};
\node[comp, minimum width=2.6cm] (user) at (0,-2.6) {Input source\\(user / network)};
\node[comp, minimum width=2.6cm] (op) at (2.2,-5.2) {Operator /\\response policy};
\draw[arr] (vendor) -- node[above, font=\scriptsize]{$\{L_i\}, r, \sigma$} (repo);
\draw[arr] (repo) -- node[right, font=\scriptsize]{load} (dram);
\draw[arr] (dram) -- node[above, font=\scriptsize]{DMA} (npu);
\draw[arr] (user) -- node[above, font=\scriptsize]{prompts} (dram);
\draw[arr] (npu) -- node[right, font=\scriptsize]{activations $\mathbf{a}_\ell$} (eg);
\draw[arr] (dram.south east) -- node[above right, font=\scriptsize, pos=0.45]{re-hash} (eg.north west);
\draw[arr] (eg) -- node[below, font=\scriptsize]{label + explanation} (op);
\node[atk] (a1) at (2.1,1.2) {A4: supply-chain\\poisoning};
\node[atk] (a2) at (1.6,-4.0) {A1: Rowhammer bit flips\\A2: cold boot / EM read};
\node[atk] (a3) at (0,-1.1) {A3/A5: crafted\\prompts, probes};
\draw[aarr] (a1) -- (repo);
\draw[aarr] (a2) -- (dram);
\draw[aarr] (a3) -- (user);
\begin{scope}[on background layer]
\node[draw=egblue, dashed, rounded corners=4pt, fit=(eg), inner sep=8pt, label={[font=\scriptsize, text=egblue]below:trusted (TEE / secure boot)}] {};
\end{scope}
\end{tikzpicture}
\caption{System model and attack surfaces. Dashed orange arrows mark where each attack class
acts. Only the \textsc{EdgeGuard} verifier, its public key and its calibration statistics are
assumed trusted.}
\label{fig:threat_model}
\end{figure}

\Cref{fig:threat_model} shows the deployment considered. A model vendor quantizes a model,
computes the Merkle root $r$ over its weight tensors, signs it with a private key and
distributes the weights with $(r,\sigma)$. On the device, the weights are loaded into DRAM and
executed by the NPU. \textsc{EdgeGuard} runs alongside. It holds the vendor's public key and
the calibration statistics, re-hashes the weights to check them against $r$, receives
intermediate activations of each inference, and reports a label and explanation to an
operator-defined response policy, such as reloading the weights, rejecting the input or
alerting.

\begin{assumption}[Trusted verifier]\label{as:trusted}
The \textsc{EdgeGuard} verification code, the public key and the calibration statistics are
protected from modification, for example by placing them in a TEE or covering them with
secure boot. The attacker cannot change them.
\end{assumption}

\begin{assumption}[Honest signing]\label{as:honest}
The weights signed by the vendor are the intended weights. A backdoor present at signing time is
outside the scope of attestation (\cref{sec:tm:goals}).
\end{assumption}

\begin{assumption}[Clean calibration]\label{as:calib}
The calibration inputs used to fit the activation monitor are free of attacks and representative
of normal use.
\end{assumption}

\begin{assumption}[Activation access]\label{as:access}
The runtime can expose per-layer activations of each monitored inference to the verifier, for
example as additional graph outputs.
\end{assumption}

% ---------------------------------------------------------------------------
\section{Adversary}
\label{sec:tm:adversary}

The adversary is assumed to have some or all of the following capabilities. Different attack
classes use different subsets.

\begin{description}[leftmargin=*,style=nextline]
  \item[C1 — DRAM write primitive.] Unprivileged code co-located on the device can induce bit
        flips in DRAM rows holding weights, via Rowhammer
        \cite{kim2014flipping,yao2020deephammer}. The attacker can aim at particular bits, subject
        to the physical constraints of the memory, and knows the model architecture and weights
        (white-box), as BFA assumes \cite{rakin2019bfa}.
  \item[C2 — DRAM read primitive.] With physical access, the attacker can read weights by a
        cold-boot attack \cite{halderman2008lest} or estimate them through an electromagnetic
        side channel, with some error.
  \item[C3 — Input control.] The attacker can submit arbitrary inputs (prompts) and observe
        outputs.
  \item[C4 — Supply-chain access.] The attacker can modify weights between signing and loading,
        for example on a mirror or in an update channel.
\end{description}

\noindent The adversary is assumed \emph{not} to be able to modify the verifier or its key
(\cref{as:trusted}), forge ECDSA signatures, or find SHA-256 collisions or second preimages.

\paragraph{Adaptivity.} A fully adaptive adversary knows the monitor's statistics and thresholds
and optimises inputs to stay under them. The evaluation in this report does \emph{not} include
such an adversary for the activation monitor. The synergy attack is a limited form of
adaptivity: the attacker tries to suppress the activation signal while modifying weights.
\Cref{sec:disc:adaptive} discusses what a fully adaptive attacker could do.

% ---------------------------------------------------------------------------
\section{Attack classes}
\label{sec:tm:attacks}

\Cref{tab:attack_classes} lists the six attack classes, the capability each requires, what it
changes, and which \textsc{EdgeGuard} signal is expected to detect it.

\begin{table}[t]
\centering
\caption{Attack classes considered, the capability they need, and the signal expected to reveal
them.}
\label{tab:attack_classes}
\small
\begin{tabularx}{\textwidth}{@{}lp{1.1cm}Xp{3.6cm}@{}}
\toprule
\textbf{Class} & \textbf{Needs} & \textbf{Changes} & \textbf{Expected signal} \\
\midrule
A1 Bit-flip attack & C1 & one or a few bits of stored weights & Merkle mismatch on the affected tensor(s) \\
A2 Weight extraction & C2 & nothing on the device (read only) & none directly; see text \\
A3 Training-data probe & C3 & nothing; inputs chosen to elicit memorised content & late-layer activation anomaly (hypothesised) \\
A4 Supply-chain poisoning & C4 & weights before loading & Merkle mismatch, often many tensors \\
A5 Adversarial input & C3 & nothing; crafted input & early/late activation anomaly \\
A6 Synergy (A1 + A5) & C1, C3 & weights and input, input tuned to mask activations & Merkle mismatch despite weak activation signal \\
\bottomrule
\end{tabularx}
\end{table}

Two classes need comment.

\paragraph{A2 is not directly detectable.} Reading memory leaves no trace in the weights or in
the activations of the device's own inferences. \textsc{EdgeGuard} cannot detect a cold-boot
read, and no claim is made that it can. The extraction simulator is included because it
quantifies how faithful a stolen copy would be (\cref{sec:eval:extraction}), and because a
stolen copy enables A3 against the original deployment or offline.

\paragraph{A3 detection rests on a hypothesis.} The extraction-probe simulator assumes that
probes produce a distinctive activation signature: elevated and low-variance late-layer
activations. If real memorisation probes do not produce such a signature, the monitor would not
detect them. Validating the hypothesis on a real model is part of the plan in
\cref{ch:conclusion}.

% ---------------------------------------------------------------------------
\section{Security goals and non-goals}
\label{sec:tm:goals}

\paragraph{Goals.}
\begin{enumerate}[label=G\arabic*,leftmargin=*]
  \item \textbf{Weight integrity.} If any byte of any attested tensor differs from the signed
        version at verification time, verification fails. Under the standard assumptions on
        SHA-256 and ECDSA, this holds with overwhelming probability.
  \item \textbf{Localisation.} When verification fails, the set of tensors reported as modified
        equals the set actually modified.
  \item \textbf{Activation anomaly detection.} Inputs whose monitored activations are atypical
        relative to calibration are flagged, at a false-positive rate the operator can choose.
  \item \textbf{Actionable output.} Each detection carries a label from a fixed set and an
        explanation naming the layers involved.
\end{enumerate}

\paragraph{Non-goals.}
\begin{itemize}[leftmargin=*]
  \item Detecting a backdoor present at signing time (\cref{as:honest}).
  \item Detecting passive reads of memory (A2).
  \item Preventing an attack, as opposed to detecting it. Recovery is left to the response policy.
  \item Protection between verifications. If weights are verified at time $t_1$ and used at
        $t_2>t_1$, a flip in between is missed until the next verification. This
        time-of-check-to-time-of-use gap is analysed in \cref{sec:disc:toctou}.
  \item Robustness of the activation monitor against a fully adaptive adversary.
\end{itemize}
